%% ==========================================================================
%%  preamble.tex --- packages, theorem environments, macros.
%%
%%  Layout deliberately mirrors Reading_3 (Barman, Krishna, Narahari,
%%  Sadhukhan, "Achieving Envy-Freeness with Limited Subsidies under
%%  Dichotomous Valuations", arXiv:2201.07419): Palatino body text, a single
%%  shared counter across all theorem-like environments, bold unpunctuated
%%  "Proof" heads, and alpha-style citation keys ([HS19], [BDN+20], ...).
%% ==========================================================================

% ---------- page geometry & fonts ----------------------------------------
\usepackage[letterpaper,margin=1.25in]{geometry}
% Palatino text + matching math, as in R3.  newpx is built on TeX Gyre
% Pagella, so it needs no URW base-35 fonts -- which a minimal TinyTeX
% installation does not ship.  If newpx is unavailable, `mathpazo' is the
% drop-in substitute wherever the URW fonts *are* present.
\usepackage[T1]{fontenc}
\usepackage[utf8]{inputenc}
\usepackage{newpxtext}
\usepackage{microtype}

% ---------- mathematics ---------------------------------------------------
% Order matters: amsmath must be loaded before newpxmath.
\usepackage{amsmath,amssymb,amsthm,mathtools}
\usepackage{newpxmath}

% ---------- floats, tables, algorithms ------------------------------------
\usepackage{booktabs}
\usepackage{graphicx}
\graphicspath{{figures/}}
\usepackage{algorithm}
\usepackage[noend]{algpseudocode}
\floatname{algorithm}{Algorithm}

% ---------- colour & links ------------------------------------------------
\usepackage{xcolor}
\definecolor{CiteGreen}{RGB}{0,120,60}
\definecolor{LinkRed}{RGB}{160,20,20}
\definecolor{UrlBlue}{RGB}{20,60,170}
\definecolor{DraftOrange}{RGB}{200,90,0}
\usepackage[
  colorlinks=true,
  citecolor=CiteGreen,
  linkcolor=LinkRed,
  urlcolor=UrlBlue,
  bookmarksnumbered=true
]{hyperref}
\usepackage{cleveref}

% ==========================================================================
%  Theorem environments.
%  R3 runs ONE counter through Definition/Theorem/Lemma/Proposition
%  (Definition 1, Definition 2, Lemma 3, Theorem 4, Proposition 5, ...).
%  Claims carry their own counter.  Keep it that way.
% ==========================================================================
\theoremstyle{plain}
\newtheorem{theorem}{Theorem}
\newtheorem{lemma}[theorem]{Lemma}
\newtheorem{proposition}[theorem]{Proposition}
\newtheorem{corollary}[theorem]{Corollary}
\newtheorem{definition}[theorem]{Definition}
\newtheorem{observation}[theorem]{Observation}
\newtheorem{conjecture}[theorem]{Conjecture}
\newtheorem{claim}{Claim}             % separate counter, as in R3

% Examples and remarks carry running argument rather than a statement, so
% they are set upright.  They keep the shared counter.
\theoremstyle{definition}
\newtheorem{example}[theorem]{Example}
\newtheorem{remark}[theorem]{Remark}
\theoremstyle{plain}

% NOTE.  cleveref cannot tell these environments apart --- they share the
% `theorem' counter, so \Cref would call every one of them "Theorem".  Refer
% to them as `Example~\ref{...}', `Remark~\ref{...}' and so on, which is also
% what R3 does.  \Cref is safe for sections.

% Bold, unpunctuated "Proof" head (amsthm defaults to italic "Proof.").
\makeatletter
\renewenvironment{proof}[1][\proofname]{\par
  \pushQED{\qed}%
  \normalfont \topsep6\p@\@plus6\p@\relax
  \trivlist
  \item[\hskip\labelsep\bfseries #1\@addpunct{}]\ignorespaces
}{%
  \popQED\endtrivlist\@endpefalse
}
\makeatother

% ==========================================================================
%  Draft notes.  Flip to \draftfalse to compile a clean copy with every
%  note suppressed --- nothing else in the document changes.
% ==========================================================================
\newif\ifdraft
\drafttrue

% ==========================================================================
%  Working-notes switch.
%
%  Two documents share this preamble:
%    main.tex     the report      -> \workingfalse (the default set here)
%    working.tex  the full draft  -> \workingtrue
%
%  Wrap in \ifworking ... \fi any sentence in a REPORT section that points at
%  material currently parked under working/.  Such a sentence then appears in
%  working.pdf, where its target exists, and vanishes from main.pdf, where it
%  would be a dangling reference.  Never put a \label inside such a guard.
% ==========================================================================
\newif\ifworking
\workingfalse

\newcommand{\draftnote}[1]{%
  \ifdraft{\color{DraftOrange}\sffamily\small[\,#1\,]}\fi}
\newcommand{\todo}[1]{\draftnote{\textbf{TODO:} #1}}
\newcommand{\verify}[1]{\draftnote{\textbf{VERIFY:} #1}}

% ==========================================================================
%  Project notation.  Fixed in glossary_fair_division_subsidies.md --- do
%  not introduce a competing symbol here; extend the glossary instead.
% ==========================================================================
\newcommand{\agents}{N}                       % agent set [n]
\newcommand{\items}{M}                        % item set, |M| = m
\newcommand{\alloc}{\mathcal{A}}              % allocation (A_1,...,A_n)
\newcommand{\allocB}{\mathcal{B}}
\newcommand{\bundle}[1]{A_{#1}}
\newcommand{\subsidy}{p}                      % subsidy vector
\newcommand{\optsubsidy}{p^{*}}               % componentwise-minimum subsidy
\newcommand{\envygraph}[1]{G_{#1}}            % envy graph
\newcommand{\edgew}[1]{w_{#1}}                % envy-graph arc weight
\newcommand{\pathw}[1]{\ell_{#1}}             % max-weight path out of a node
\newcommand{\SW}{\mathrm{SW}}                 % utilitarian welfare
\newcommand{\uspread}{\mathrm{uspread}}       % spread of the total-cost function
\newcommand{\paidsets}{\mathcal{P}}           % admissible paid sets of a state
\newcommand{\Mset}{\mathrm{M}}                % set of maximally subsidised agents

% Marginals.  Delta^+ is the gain from adding g to S; Delta^- the loss from
% removing g from S.  Distinct objects once v_i is not additive.
\newcommand{\margplus}[3]{\Delta^{+}_{#1}(#2,#3)}
\newcommand{\margminus}[3]{\Delta^{-}_{#1}(#2,#3)}

% Cost (disutility) form of a negative dichotomous valuation: c_i := -v_i.
\newcommand{\cost}{c}

% Peel frame (working/approach_3.tex).  A workload profile W = (W_1,...,W_n)
% records which chores each agent is still on the hook for.  Not yet in the
% glossary --- promote it there before this notation leaves working/.
\newcommand{\workprofile}{W}

\newcommand{\N}{\mathbb{N}}
\newcommand{\R}{\mathbb{R}}
\newcommand{\Z}{\mathbb{Z}}
\newcommand{\set}[1]{\left\{#1\right\}}
\newcommand{\abs}[1]{\left\lvert#1\right\rvert}
\DeclareMathOperator*{\argmax}{arg\,max}
\DeclareMathOperator*{\argmin}{arg\,min}
